\chapter{Hausdorffabstand und kompakte Mengen}
\label{b47guide:ch:hausdorff}

In diesem Kapitel ist $(X,d)$ ein nichtleerer metrischer Raum.
Der Hausdorffabstand macht seine nichtleeren kompakten Teilmengen
selbst zu einem metrischen Raum.

\section{Abstand zwischen kompakten Mengen}

\begin{definition}[Hausdorffabstand]
\label{b47guide:def:hausdorff}
Es sei $\mathcal K(X)$ die Menge aller nichtleeren kompakten Teilmengen
von $X$. Für $A,B\in\mathcal K(X)$ heißt
\[
 d_H(A,B)=\max\left\{
   \sup_{a\in A}d(a,B),\ \sup_{b\in B}d(b,A)
 \right\}
\]
der \emph{Hausdorffabstand} von $A$ und $B$. Die Abstände zu Teilmengen
sind die bereits definierten Punkt-Mengen-Abstände.
\end{definition}

\begin{theorem}[Hausdorffmetrik und Einbettung der Punkte]
\label{b47guide:thm:hausdorff-metric}
$d_H$ ist eine endliche Metrik auf $\mathcal K(X)$.
Für alle $x,y\in X$ gilt
\[
 d_H(\{x\},\{y\})=d(x,y).
\]
Die Abbildung $x\mapsto\{x\}$ ist daher eine isometrische Einbettung
von $X$ in $\mathcal K(X)$. \cite{hausdorff-notes}
\end{theorem}

\begin{theorem}[Abstandsfunktionen und Durchmesser]
\label{b47guide:thm:hausdorff-distances}
Für alle $A,B\in\mathcal K(X)$ gelten
\[
 d_H(A,B)=\sup_{x\in X}|d(x,A)-d(x,B)|,
 \qquad
 |\operatorname{diam}A-\operatorname{diam}B|\leq2d_H(A,B).
\]
Die erste Gleichheit gilt auch bei unbeschränktem $X$.
Somit bedeutet $d_H(A_n,A)\to0$ genau die gleichmäßige Konvergenz
der Abstandsfunktionen $x\mapsto d(x,A_n)$ gegen $x\mapsto d(x,A)$
auf ganz $X$. \cite{hausdorff-distance}
\end{theorem}

\section{Vollständigkeit und Auswahl kompakter Mengen}

\begin{theorem}[Vollständigkeit des Raums kompakter Mengen]
\label{b47guide:thm:hausdorff-complete}
Der metrische Raum $(\mathcal K(X),d_H)$ ist genau dann vollständig,
wenn $(X,d)$ vollständig ist.
Insbesondere besitzt bei vollständigem $X$ jede Cauchy-Folge
nichtleerer kompakter Mengen einen eindeutig bestimmten,
nichtleeren kompakten Grenzwert bezüglich $d_H$.
\cite{hausdorff-notes}
\end{theorem}

\begin{theorem}[Endliche Approximation, totale Beschränktheit und Separabilität]
\label{b47guide:thm:hausdorff-approximation}
Die nichtleeren endlichen Teilmengen von $X$ liegen dicht in
$\mathcal K(X)$. Ist $D\subseteq X$ dicht, so liegen bereits die
nichtleeren endlichen Teilmengen von $D$ dicht in $\mathcal K(X)$.
Ferner gelten die Äquivalenzen
\[
\begin{aligned}
 \mathcal K(X)\text{ total beschränkt}
   &\quad\Longleftrightarrow\quad X\text{ total beschränkt},\\
 \mathcal K(X)\text{ separabel}
   &\quad\Longleftrightarrow\quad X\text{ separabel}.
\end{aligned}
\]
\end{theorem}

\begin{theorem}[Blaschkes Auswahlsatz in metrischer Form]
\label{b47guide:thm:blaschke}
$(\mathcal K(X),d_H)$ ist genau dann kompakt, wenn $X$ kompakt ist.
Insbesondere besitzt bei kompaktem $X$ jede Folge $(K_n)$ nichtleerer
kompakter Teilmengen eine Teilfolge $(K_{n_j})$ und eine Menge
$K\in\mathcal K(X)$ mit
\[
 d_H(K_{n_j},K)\longrightarrow0.
\]
\cite{hausdorff-notes}
\end{theorem}

\section{Abstandskontrolle für Bilder}

\begin{theorem}[Lipschitz-Abbildungen auf Räumen kompakter Mengen]
\label{b47guide:thm:hausdorff-images}
Sei $(Y,d_Y)$ ein metrischer Raum und sei $f:X\to Y$ eine
$L$-Lipschitz-Abbildung, $L\geq0$. Dann ist die Abbildung
\[
 f_*:\mathcal K(X)\longrightarrow\mathcal K(Y),\qquad A\longmapsto f[A],
\]
bezüglich der jeweiligen Hausdorffmetriken ebenfalls $L$-Lipschitz:
\[
 d_H^{Y}(f[A],f[B])\leq L\,d_H^{X}(A,B)
 \qquad(A,B\in\mathcal K(X)).
\]
\end{theorem}